\subsection{Class-four coefficient deformation and private point frames}

% Mathematical recovery transcription; see recovery provenance below.
\paragraph{Recovery provenance.}
This is a complete mathematical transcription of the three-part E.217 proof delivered in conversation during the filesystem outage. It is not a byte-identical recovery of a saved file: no E.217 file, compilation, or executable certificate existed during that outage. Root read all three parts, and ghost\_audit and literature completed positive full independent proof reviews. This version includes literature's two requested precision additions: generator values in the cohomology calculation are Lie cocycle values under BCH, and the higher literal B-cocycle is aligned before invoking the fixed-B tangent torsor at a new prime.

\subsubsection*{Statement and retained data}
Let $N_4$ be the actual geometrically pro-five free nilpotent group of class four for $\mathbb P^1-\{0,1,\infty\}$, with its actual tangential arithmetic action. Since $4<5$, integral Lazard coordinates are available. Fix $n\ge1$, an actual global section $s_n$ in its coefficient-$5^n$ quotient, and a continuous global higher B-coordinate Kummer cocycle $B_{n+1}$ reducing literally to that of $s_n$.

Let $S$ contain $5,\infty$, all ramification of the retained section, and all ramification of $B_{n+1}$. Supply compatible next-coefficient local lifts at every old place in $S$. At $83$ retain an actual A-point with A-valuation a five-adic unit, and at $173$ retain an actual B-point with B-valuation a five-adic unit. These are the private point conditions used below. The local lifts retain the literal lower section and the literal higher B-cocycle, but all actual finite lower aligning frames are allowed.

\begin{result}{E.217}{Proved}{Class-four coefficient interpolation with private point frames}
Under these hypotheses, a global next-coefficient class-four lift exists, retaining the lower global section and the prescribed higher B-cocycle and matching all the supplied old local classes. New ramification can be confined to finitely many arbitrarily large primes $q$ satisfying
\[
 \operatorname{ord}(q\bmod5)=4,
 \qquad v_5(q^4-1)=1.
 \tag{1}
\]
At each new prime its local class is that of an actual A-point $q^{m_q}$, with a positive integer $m_q$ divisible by $5^n$ and by any specified integer prime to five.

One may retain a finite arithmetic field with commutator of order prime to five and a specified compatible Frobenius carrying additional CRT, radical-root, or other-coefficient data. In that version the exact remaining arithmetic characters are retained; the theorem requires their pairing with the old first-coordinate discrepancy to vanish. A sufficient condition is that the old first-coordinate discrepancy be the restriction of one global abelian tangent class, unramified outside the old set.

Finite iteration is allowed only with compatible supplied old local lifts through the requested precision, for example reductions of fixed integral local sections. This theorem is a coefficient refinement in fixed class four. It does not prove a lift from class three to class four, a cofinal ordinary-preserving embedding theorem, or an infinite branch.
\end{result}

\begin{proof}
\par\medskip\noindent\textbf{1. The fixed-B tangent module and the actual action}\par\nobreak\smallskip
Write
\[
 C=[X,Y],\quad U=[X,C],\quad V=[Y,C],
\]
\[
 A=[X,U],\qquad B=[Y,U]=[X,V],\qquad W=[Y,V].
\]
The fixed-B tangent ideal has basis
\[
 I=\langle X,C,U,V,A,B,W\rangle_{\F_5}.
\]
Let $D=I^*(1)$. Its graded dual weights and multiplicities are
\[
 0\ (1),\qquad 3\ (1),\qquad 2\ (2),\qquad 1\ (3)
 \quad\text{modulo }4.
 \tag{2}
\]
The degree-two actual path coefficient is $(\chi^2-1)C/24$ and is removed by a fixed integral peripheral change of frame, or may equivalently be retained in every formula. The degree-three arithmetic IA action fixes $X$ and sends $Y$ into weight four. It acts trivially on $I$: it fixes $X$, and its effect on every bracket of weight at least two has weight at least five. Thus the minimal action on $I$ is
\[
 J=P_3\rtimes C_4,
 \tag{3}
\]
where $P_3$ is the free class-three exponent-five group. The degree-three arithmetic action is nevertheless retained in the actual section and cusp Frobenius data. Its being trivial on this one tangent ideal does not license deleting it from the arithmetic extension.

The two private inertia groups give a Frattini basis of the geometric image. They are unramified in the based arithmetic action, which is unramified away from five. Consequently the actual joint representation retains the full finite geometric class-four quotient above its arithmetic action field. The kernel of its action on $I$ is precisely the central weight-four subgroup, giving (3).

\par\medskip\noindent\textbf{2. Minimal-action cohomology and cyclic detection}\par\nobreak\smallskip
First compute on the free class-four exponent-five group $P_4$. The cocycle graph for coefficients $D$ still has nilpotence class at most four, so BCH identifies its group cocycles with the corresponding Lie cocycles. After the coprime $C_4$ normalization, the only possible generator values are
\[
 f(X)=\alpha,\qquad f(Y)=\beta,
 \qquad \alpha,\beta\in L_4^*.
 \tag{4}
\]
These are \emph{Lie} cocycle values. The group cocycle at $\exp X$ or $\exp Y$ can have determined lower-dual components; its top components remain respectively $\alpha$ and $\beta$, which are all that the detection argument uses. There are no invariant coboundaries: the only weight-zero dual line is $X^*$, and it is geometrically fixed.

Writing $\alpha(A),\alpha(B),\alpha(W)$ for its three coordinates, the Lie cocycle identity gives
\[
 f(A)(X)=-\alpha(A),\quad
 f(B)(X)=-\alpha(B),\quad
 f(W)(X)=-\alpha(W).
 \tag{5}
\]
For example $f(C)=X\beta-Y\alpha$ and the next two brackets give (5) directly. The degree-four relations defining $P_3$ therefore force $\alpha=0$ and leave $\beta$ arbitrary. Thus
\[
 H^1(J,D)\simeq\F_5^3.
 \tag{6}
\]
We do not replace (6) by an erroneous vanishing assertion.

Every nonzero class in (6) has a nonzero top component on the cyclic subgroup generated by $Y$. Such a component cannot be a cyclic coboundary, because $(\operatorname{Ad}Y-1)D$ has zero top-dual component. Hence
\[
 \Gamma(J,D):=
 \ker\!\left(H^1(J,D)\longrightarrow
 \prod_{C\le J\ \mathrm{cyclic}}H^1(C,D)\right)=0.
 \tag{7}
\]

\par\medskip\noindent\textbf{3. A weak global reference lift}\par\nobreak\smallskip
The next-coefficient obstruction lies in $H^2(G_{\Q,S},I)$. It is locally zero at old places by hypothesis. Outside $S$, lift the unramified Frobenius element and set inertia to one. This also respects the prescribed higher B-cocycle because its ramification was included in $S$.

Equation (7), together with Chebotarev on the kernel of the actual finite action, detects every nonzero globally locally-trivial candidate in the dual degree-one group. Choose finitely many temporary places $T_0$ detecting the finite old dual space. A dual class over $G_{\Q,S\cup T_0}$ that is locally zero at those temporary places has no new inertia and descends to $G_{\Q,S}$, where it is detected. Thus the degree-one local-global kernel vanishes after adding $T_0$. Finite Poitou--Tate duality kills the dual degree-two kernel. The locally soluble obstruction vanishes and gives a weak global next lift. We subsequently prescribe unramified targets at $T_0$ and remove its temporary ramification.

\par\medskip\noindent\textbf{4. The finite central frames at old private places}\par\nobreak\smallskip
At the A-private point the geometric inertia is conjugate to $\exp(mX)$ with $m\in\Z_5^\times$. The inertia invariants in $I$ are $\F_5X\oplus L_4$. At the B-private point they are $L_4$. Frobenius on the first summand has weight one, and on $L_4$ it has weight four. Thus the unramified tangent classes at either private point are exactly the three weight-four directions.

These are genuine available changes of the finite lower aligning frame. Let
\[
 e=v_5(p^4-1),\qquad n\ge e,
\]
and take an integral $V_4\in L_4$. The geometric element
\[
 h=\exp(5^{\,n-e}V_4)
 \tag{8}
\]
stabilizes the lower finite section. It commutes with inertia \emph{exactly}, since weight four is central. Its lifted Frobenius boundary at the next coefficient layer is the unit multiple
\[
 \pm\frac{p^4-1}{5^e}\,V_4.
 \tag{9}
\]
Hence every unramified weight-four tangent direction occurs. This is valid already at $n=e=1$; the frame need not reduce to the identity modulo five, only stabilize the retained lower section. For $p=83,173$, the selected primes have $e=1$.

The local Tate edge pairing of an unramified $I$-class with a dual class evaluates its dual inertia residue on $L_4$. On the A-private inertia this top residue is $m\alpha$; on the B-private inertia it is $m\beta$. Orthogonality to the available frames therefore forces
\[
 \alpha=\beta=0.
 \tag{10}
\]
If an extra retained arithmetic field ramifies at one of these places, its inertia component acts trivially on the actual geometric group and on $I$. Commuting that component with the full geometric image forces its cocycle value into $D^{P_4}=\F_5X^*$, which evaluates to zero on $L_4$. It therefore cannot cancel the top pairing in (10).

\par\medskip\noindent\textbf{5. New ordinary primes, including literal B and Frobenius}\par\nobreak\smallskip
At a new prime satisfying (1), require the retained lower Frobenius to be geometrically conjugate to the actual cusp-zero graph. This is an explicit finite Frobenius condition. It is not inferred merely from unramifiedness in class four.

At such a prime the fixed-B tangent cohomology consists of
\[
 H^1_{\mathrm{nr}}(\Q_q,I)=L_4,
 \qquad\text{and one tame A-line }\F_5X.
 \tag{11}
\]
Indeed weights one through three have $q^i-1$ units, weight four is the unramified resonance, and in $I(-1)$ only the A-direction is Frobenius invariant. An actual point $q^{5^n a}$ realizes the tame line. The higher prescribed B-cocycle is unramified and has zero local Kummer class because $q-1$ is a unit. First change the point representative by a congruence Y-alignment so that its higher B-cocycle is \emph{literally} the prescribed one; this preserves the lower cusp frame. Only then compare in the $I$-torsor.

All differences in the first term of (11) are absorbed by the central frames (8)--(9), with $e=1$. It follows that every next-coefficient fixed-B local lift above the specified lower cusp class is represented by an actual ordinary A-point, up to the allowed finite frame. This includes Frobenius, not just inertia. Choose the residue $a$ by its positive integer valuation representative and impose any desired prime-to-five divisibility by CRT.

Orthogonality to the full unramified term in (11) forces every new dual tame residue to vanish. Thus enlarging the auxiliary set creates no new allowed dual classes, even though the raw local dual module has weight-four residue invariants. This is the point missed by an argument that freezes an arbitrary lower finite frame.

\par\medskip\noindent\textbf{6. Detection in the full retained representation}\par\nobreak\smallskip
Let $A_{\mathrm{ar}}$ carry the actual based class-four action, cusp data, and $\chi$ modulo $25$. It is unramified outside five. The joint image of the retained global section is
\[
 P_{4,n}\rtimes A_{\mathrm{ar}},
 \tag{12}
\]
with full finite geometric group, by private inertia and Frattini. Choose an actual arithmetic element whose cyclotomic residue is primitive and whose Fermat digit is nonzero. Permitted auxiliary Frobenii are the geometric conjugates of its cusp graph. Their cyclotomic component gives (1).

Let a global dual class satisfy the private orthogonality conditions and pair trivially with all permitted ordinary auxiliary targets. Let $N$ be the kernel of the retained joint representation (12), viewed inside the finite extension that also carries the dual cocycle. Thus $N$ is not defined by killing the dual cocycle itself. Comparing permitted Frobenii $g$ and $gn$ for $n\in N$ shows that the value on $n$ annihilates every conjugate of the peripheral tangent $X$.

These conjugates span $I$. For example the values of $\operatorname{Ad}(\exp(tY))X$ for $t=0,1,2,3$ give $X,C,V,W$ by an invertible Vandermonde system. Conjugation by $\exp(sX)$ then gives $U,A$, and conjugating $V$ gives $B$. The denominators are at most $3!$, all units. Hence the class vanishes on $N$ and descends to (12).

Its geometric restriction factors through the exponent-five group $P_4$. On $D$, the adjoint nilpotence degree is at most three, so
\[
 1+g+g^2+g^3+g^4=0
\]
for the relevant geometric actions. The cocycle therefore kills fifth powers and their normal closure. After the coprime scalar normalization its geometric restriction has the six parameters (4). Additional actual arithmetic IA directions may impose further relations, but they cannot create extra parameters. The two private conditions (10) kill all six. Inflation--restriction then leaves only arithmetic classes with values in
\[
 D^{P_{4,n}}=\F_5X^*.
\]
For the actual based arithmetic field alone, its five-abelian quotient is unramified outside five and is cyclotomic. The nonzero Fermat digit of the permitted Frobenius detects its unique possible character. Thus no dual obstruction remains.

Only finitely many old dual classes occur, so finitely many permitted primes detect them. By the preceding no-new-residue argument every dual class over the enlarged ramification set descends to the old set. Finite Poitou--Tate gives the desired tangent correction at old places, modulo the available finite private frames. At new places every correction is ordinary by (11). Prescribe unramified targets at $T_0$ to remove the temporary ramification of the weak reference lift.

\par\medskip\noindent\textbf{7. Additional retained arithmetic and reciprocity}\par\nobreak\smallskip
Let $F$ be a finite retained field with commutator of order prime to five. The geometric group in (12) lies in the derived subgroup, so its intersection with $F$ is arithmetic. Full geometric image therefore survives above every compatible arithmetic element of the compositum. Specify an \emph{actual} compatible element $\sigma$ carrying the desired CRT, radical-root, and other-coefficient prescriptions, with primitive cyclotomic residue and nonzero Fermat digit. Compatibility of those prescriptions must be proved before applying this theorem.

The same kernel and geometric arguments leave exactly the arithmetic characters
\[
 \psi\in\operatorname{Hom}(A',\F_5),\qquad \psi(\sigma)=0,
 \tag{13}
\]
where $A'$ is the actual retained arithmetic compositum. Their local pairing depends only on the first-coordinate tangent projection. If the old projection is the restriction of one global Kummer tangent class, its pairing with each character in (13) sums to zero by global reciprocity. Outside the old set the two classes are unramified, and at the permitted new primes the character vanishes on $\sigma$. Thus the supplied compatibility annihilates precisely the remaining dual kernel, and Poitou--Tate supplies the correction.

One can obtain a baseline using the unrestricted temporary-prime construction and then apply this retained-field argument, including the temporary places among the old places with compatible unramified targets. This does not assert arbitrary interpolation after arbitrary additional arithmetic restrictions.

\par\medskip\noindent\textbf{8. Finite iteration and the class-jump boundary}\par\nobreak\smallskip
Freeze a positive integer valuation at each newly introduced ordinary point and retain that actual point in later coefficient refinements. Original old local targets must likewise have supplied compatible lifts through the required precision, for example come from fixed integral local sections. This gives finite iteration. It does not assert that arbitrary old finite local classes lift indefinitely.

For primitive/minimal new primes, the lower unramified cusp condition has a fixed finite test. Weights one through three have unique local aligning classes, and the remaining central cokernel is
\[
 (\Z/5^n(4))^3/(q^4-1)(\Z/5^n(4))^3\simeq\F_5^3.
\]
Reduction modulo five is an isomorphism on this cokernel. Thus the actual finite class-four cusp graph modulo five detects the required lower unramified matching condition at every coefficient precision. This helps retain an already admissible class-four predecessor; it does not construct that predecessor from a class-three section.

Finally, the maps $N_4(\Z/5^{n+1})\to N_4(\Z/5^n)$ and $N_4(\Z/5^n)\to N_3(\Z/5^n)$ are different embedding problems. A mismatch already visible in a retained finite class-four quotient cannot be repaired by later coefficient digits. No class-three-to-four jump, infinite ordinary completed pair, or nonempty all-depth catalog follows from this theorem.
\end{proof}

\paragraph{Source scope.}
The local Tate pairing and finite Poitou--Tate exactness are those of Milne, \emph{Arithmetic Duality Theorems}, Chapter I, Section 2 and Theorems 4.9--4.10, pp.~56--57, independently opened earlier in the run. The actual arithmetic path coefficient and its degree-three direction are the established B.23 inputs. All finite-module, central-frame, and local-inertia calculations used here are proved above. There was no new executable certificate for E.217.
